\section{Method}

\subsection{Blind-Spot Coverage as a Verifiable Reward}

Given a query $q$, let $G=\{g_1,\dots,g_m\}$ be a gold
evaluation-dimension set, and let $\hat{R}=\{r_1,\dots,r_n\}$ be a generated
criteria set. In paper-facing hard-gold evaluation, $G$ denotes human-gold
evaluation dimensions; in training-time reward computation, the hidden
criteria are proxy-gold targets from disjoint query pools. These roles are not
interchangeable. We
compute semantic similarity between gold and generated dimensions and define
coverage as
\[
\cov(\gold,\pred) = \frac{1}{|\gold|}\sum_i
\mathbb{1}\left[\max_j \cos(e(g_i), e(r_j)) \ge \tau\right].
\]
The blind-spot rate is $\blind = 1-\cov$. This metric differs from text
similarity to a reference answer: it asks whether every human evaluation
dimension is semantically represented somewhere in the generated set, allowing
lexical variation while preserving dimension-level accountability.

BSC becomes a reward only after adding two controls. First, we penalize
redundancy within the generated set:
\[
\red(\pred) =
\frac{|\{(r_a,r_b): a<b, \cos(e(r_a), e(r_b)) \ge \tau'\}|}{n},
\]
This prevents reward hacking through repeated paraphrases of the same
dimension. Second, a verifier estimates whether each generated criterion is
atomic, decidable, relevant, and non-hallucinated. We define
$\hall(\pred)=1-R_{\mathrm{valid}}$, where $R_{\mathrm{valid}}$ is the
fraction of generated criteria passing all four checks. Invalid,
undecidable, irrelevant, hallucinated, or otherwise unverifiable criteria
receive no validity credit. The resulting scalar reward is
\[
R_{\mathrm{BSC}} =
w_{\mathrm{cov}}\cov(\gold,\pred)
+ w_{\mathrm{valid}}R_{\mathrm{valid}}(\pred,q)
- w_{\mathrm{red}}\red(\pred).
\]
The default implementation uses $w_{\mathrm{cov}}=1.0$,
$w_{\mathrm{valid}}=0.5$, and $w_{\mathrm{red}}=0.5$, with
$\tau=0.75$ and $\tau'=0.85$. Unless otherwise stated, semantic similarities
are computed with the fixed \texttt{BAAI/bge-large-en-v1.5} embedding model
from the BGE family~\cite{xiao2023cpack}; toy smoke tests use a separate
token-overlap embedder and are never used for paper-facing BSC claims.

The embedding model, thresholds, and reward weights are fixed by protocol
before final hard-gold evaluation rather than tuned on
RubricBench \texttt{test\_main}. Absolute BSC values are accompanied by
threshold sweeps, matched/unmatched human-audit packs, and bootstrap confidence
intervals. Method comparisons are reportable only when the compared rows use
the same embedding model, thresholds, verifier-backed validity source, and
paired bootstrap protocol.

The key point is that verifiability comes from dimension-level semantic
coverage, not from matching a single reference string. This lets us apply
an RLVR-style optimization test to an open-ended task: many different criteria
wordings can receive credit if they cover the same human evaluation dimension,
while redundant or unverifiable criteria are penalized. BSC therefore supplies a
checkable reward for subjective criteria elicitation without reducing the task
to exact-answer prediction.
This is not prompt tuning with a downstream preference signal: the reward is
computed from hidden evaluation dimensions under the relevant evidence role
(human-gold for hard-gold evaluation and proxy-gold during training),
intra-output redundancy, and fail-closed validity checks over atomicity,
decidability, relevance, and non-hallucination, so a policy cannot receive
paper-facing credit merely by changing surface wording.

\subsection{Blind-Spot Attribution}

BSC identifies missed gold dimensions but does not by itself explain what kind
of evaluation signal is missing. We therefore construct a Blind-Spot Map by
attributing each uncovered human-gold dimension to a fixed category such as
factuality, completeness, constraint following, safety, domain knowledge,
evidence grounding, or intent and reasoning. For a query $q$, the blind-spot set is
\[
B_q=\left\{g_i \in G: \max_j \cos(e(g_i), e(r_j)) < \tau\right\}.
\]
Let $c_i = C(g_i)$ be the category assigned to human-gold dimension $g_i$. The
category-specific blind-spot rate is
\[
\blind_k(q) =
\frac{|\{g_i \in B_q: c_i=k\}|}{|\{g_i \in G: c_i=k\}|},
\]
computed only for categories present in the query. This yields an auditable
model-by-category map of where generated criteria miss human evaluation
dimensions.

\subsection{Category-Balanced Reward}

Some criteria categories are more frequent than others. Optimizing raw coverage
alone may cover common dimensions while leaving rare categories unresolved.
When gold category labels are available, we add a query-local
category-balanced coverage term:
\[
R_{\mathrm{bal\_cov}} =
\frac{1}{|K_q|}\sum_{k \in K_q}
\frac{|\{g_i \in \gold: C(g_i)=k,\ \max_j \cos(e(g_i), e(r_j)) \ge \tau\}|}
{|\{g_i \in \gold: C(g_i)=k\}|}.
\]
The balanced reward is
\[
R =
0.7R_{\mathrm{cov}} + 0.3R_{\mathrm{bal\_cov}}
+ 0.5R_{\mathrm{valid}} - 0.5R_{\mathrm{red}}.
\]
This turns the blind-spot map into a training signal: the policy is rewarded
for covering long-tail human dimensions, not merely for increasing average
criteria length.

\subsection{Fail-Closed Meta-Verification}

Proxy-gold criteria are useful only if low-quality generated criteria are
prevented from entering training. BlindSpot-RL therefore uses a fail-closed
meta-verifier. The rule verifier rejects empty, too-short, generic,
overly-long, or meta-instruction-like criteria. The optional API verifier
returns structured judgments for atomicity, decidability, relevance, and
non-hallucination. Items that cannot be parsed, cannot be judged, or trigger
verifier errors are marked invalid rather than silently accepted; undecidable
criteria therefore reduce $R_{\mathrm{valid}}$ just like hallucinated or
irrelevant criteria. A
domain-aware verifier policy relaxes relevance for source-specific evaluation
needs, such as medical risk disclosure for HealthBench, search evidence and
freshness for BEIR/NQ, or instruction constraints for IFBench.

\subsection{Two-Stage Training}

We use a two-stage training recipe. Stage 1 performs supervised fine-tuning on
criteria elicited by a multi-teacher union pipeline. Teacher outputs are grouped
by query, merged, semantically deduplicated, verifier-filtered, and converted
to proxy-gold SFT records. Stage 2 applies GRPO/RLVR using the BSC reward.
The verl-compatible reward hook reads hidden gold criteria from
\texttt{ground\_truth} or \texttt{extra\_info}, parses the model response into
atomic criteria, applies the fail-closed rule verifier for the validity term
unless response-aligned verifier flags are supplied, and returns the BSC reward.
Empty or unparsable outputs receive a negative reward, discouraging format
shortcuts.
The hidden criteria used by this reward during training are proxy-gold training
targets, not RubricBench \texttt{test\_main} human-gold dimensions. This
separation is essential to the research claim: BSC makes an open-ended
criteria-elicitation task rewardable, while the hard-gold split remains an
external test of whether the learned policy covers human evaluation dimensions
it was not trained on. Thus the RLVR transfer is evaluated as a protocol for
semantic reward optimization under holdout evidence, rather than as direct
optimization on the final evaluation set.

\subsection{Contamination-Aware Data Protocol}

The method relies on strict separation between evidence roles. Human-gold
RubricBench is physically split into \texttt{train\_seed} (500 examples),
\texttt{dev} (150), and \texttt{test\_main} (497 examples; 491 unique
queries). Only \texttt{test\_main} is
allowed in the main BSC evaluation; it is not used to choose prompts, tune
reward weights, calibrate the verifier, select checkpoints, or construct
proxy-gold targets. ResearchRubrics is split into 60
\texttt{train\_seed} examples and 41 \texttt{dev\_test} examples for
generalization analysis. RewardBench is first split by query into
\texttt{sft\_proxy\_train}, \texttt{dev}, and
\texttt{downstream\_holdout} partitions (1,791/597/597 records); after
hard-gold and downstream holdout filtering, the clean proxy-training partition
contains 1,785 records. RewardBench-2 and JudgeBench are full downstream
holdouts. The current downstream holdout audits cover 597 RewardBench records
with 551 normalized unique queries, 620 JudgeBench records with 528 normalized
unique queries, and 1,815 RewardBench-2 records with 1,814 raw stripped unique
queries or 1,813 normalized unique queries. IFBench, WritingBench, HealthBench,
and BEIR/NQ are query pools that
may enter training only after multi-teacher criteria elicitation and verifier
filtering. Before any SFT or RL step, we run query-level holdout filtering
against RubricBench \texttt{test\_main} and the downstream RewardBench,
JudgeBench, and RewardBench-2 holdouts; any proxy-training example with an exact
normalized query match is removed. We then run contamination audits against
\texttt{test\_main} and each downstream holdout, and every audit must report
\texttt{overlap\_query\_count == 0}. This zero-overlap count is not sufficient
by itself: if any configured training-side artifact is missing, the audit is
marked \texttt{artifact\_status=blocked} and
\texttt{overlap\_status=not\_auditable}, so it cannot support C0. The filtering reports record SHA256
digests for the holdout, input, and output files, and the final audit reports
record SHA256 digests for the holdout and each training-side artifact. The
claim-evidence matrix checks that each filtering stage's output digest matches
the next stage's input digest, and that the final clean proxy-train digest is
the same artifact inspected by the hard-gold and downstream audits. It also
checks that the verifier-filtering report binds the raw teacher criteria, the
filtered teacher criteria, the Meta-Verifier provider, the SFT-data preflight
report, and the meta-verifier budget report by SHA256, and that the filtered
teacher digest is the same input digest consumed by the proxy-gold construction
report. Finally, it checks that the SFT JSONL, proxy-gold JSONL, and GRPO
parquet digests recorded by the data-construction reports match the artifacts
inspected by every final holdout audit. These
audits inspect both visible prompt
fields and nested reward metadata such as \texttt{extra\_info.query}, so
\texttt{proxy\_gold\_verl.parquet} is checked by its original query rather than
only by the formatted training prompt. This protocol is designed to make the
main and downstream utility claims resistant to training/evaluation
contamination concerns.

\subsection{Evidence-Gated Claims}

All paper claims are linked to artifacts through a claim-evidence matrix. A
claim is draftable only if its required files exist and its metric gates pass.
This prevents the manuscript from overstating toy or incomplete evidence.
